\section{Introduction}\label{introduction}
Unknown-environment exploration must discover accessible space while acquiring enough views to reconstruct its surfaces. These objectives overlap but are not interchangeable. A robot can map a corridor while missing the back of a cabinet or the interior of a shelf; inspecting those surfaces can also delay exploration elsewhere. The practical question is how category information should affect observation allocation when both the map and object structure remain incomplete.

We study a joint unknown-coverage and measured-surface planner in which category votes weight closed and open structural alternatives. Candidate routes come from acquired occupancy and surface data, and each executed motion is followed by a new observation. The measured map remains separate from hypothetical completion. A neutral class-prior control G and a category-conditioned policy S share instance associations, geometry, routes, mapping and execution at any identical observed state. This comparison isolates the class-prior contribution within the declared frontend; it does not compare semantic perception against a wholly geometric sensing system.

A recurring failure is that available category information produces no target change. We make this behavior explicit: the semantic score difference is the class-probability difference multiplied by each route's difference in predicted structural observation opportunity. If geometric evidence is overconfident, the two structural predictions expose similar area, or the score change does not overcome the ranking margin, G and S choose the same route. A target change is only an intermediate result; it must improve actual coverage or reconstruction to support a task-level advantage.

Our evaluation connects forecast scores to measured route outcomes and complete unknown exploration. Four new layouts span chain and cycle room topologies; paired visibility/acquisition policies share the fitted full posterior, original quadrature, and execution rules. Twenty-four fixed branches connect both forecasts to the same acquired outcomes. Independent procedural objects test extent and structural probability, while earlier G/S floorplan permutations and native correct/neutral/swapped-cue branches isolate narrower prior effects. Supplied-graph studies and physical integration pilots remain supporting evidence with separate assumptions.

The contributions are threefold. First, an implementation-traced joint planning formulation separates acquired mapping, hypothetical structural reward and realized execution outcomes. Second, an exact score decomposition explains why semantic probabilities can change while the selected target remains fixed, and motivates independently testing structural evidence scale. Third, a matched evaluation connects probability changes to native candidate choices and complete exploration, retaining failures, ties and adverse tradeoffs. The current evidence establishes exploration feasibility and improved structural probability loss on a held-out procedural set. It does not establish a general semantic advantage in unknown-space coverage or reconstruction.

\section{Related Work}\label{related-work}
Active mapping selects motion to improve an environmental representation while accounting for estimation~\cite{ref10}. Hierarchical systems such as TARE~\cite{ref5} and Active Neural SLAM~\cite{ref1} separate strategic exploration from local execution. Our virtual study uses the existing observed-state route generator and primitive executor to examine the additional value of a category-conditioned surface reward. Exact poses remove localization uncertainty from that comparison; they do not validate physical SLAM accuracy.

Reconstruction-driven next-best-view planning optimizes surface acquisition; GenNBV~\cite{ref6} learns a generalizable policy in a larger viewpoint space. Semantic exploration and inspection methods use category or scene structure to prioritize observations~\cite{ref2,ref7,ref3,ref8}. Here the comparison concerns two explicit structural alternatives on a ground robot's available routes. Measured surfaces are integrated volumetrically~\cite{ref4,ref11}, while predicted geometry affects only observation selection. Adapted mechanisms and historical native-system integration results in the supplement are not rankings of those complete original systems.

Belief-based planning and robotic information gathering provide the decision background~\cite{ref12,ref13}. The present loss-based shape weights follow a generalized updating approach~\cite{bissiri2016}; the structural-temperature adjustment uses standard probability calibration~\cite{pmlr-v70-guo17a}. Conditional depth CRPS, class NLL and interval coverage assess different properties. We therefore evaluate them separately and treat a standard split-conformal interval correction~\cite{Angelopoulos_2023} as a diagnostic, not a new navigation algorithm or a guarantee of improved reconstruction.
